\appendix
\section*{Приложение А. Структура проекта}
\addcontentsline{toc}{section}{Приложение А. Структура проекта}

Ниже приведена файловая структура проекта (служебные и сгенерированные каталоги опущены).
Переиспользуемый код собран в пакете \code{src/} и разбит по доменам; запускаемые сценарии вынесены
в \code{scripts/}, а гиперпараметры — в \code{configs/}.

{\footnotesize
\begin{verbatim}
coursework/
|-- src/                      переиспользуемый код (пакеты по доменам)
|   |-- common/               конфигурация, пути, ввод-вывод, кэш LLM
|   |-- data/                 загрузка OpenAlex, сборка истории авторов
|   |-- embeddings/           векторизация статей и метадокументов (e5-base)
|   |-- clustering/
|   |   |-- paper_topics/     UMAP, HDBSCAN, бейзлайны, c-TF-IDF, MMR
|   |   \-- meta_topics/      метакластеры и UI-маппинги
|   |-- recommendation/       модель автора, обучение, retrieval-пайплайн
|   |-- llm/                  разметка тем, LLM-оценка рекомендаций
|   |-- evaluation/           метрики кластеризации, retrieval, покрытия
|   \-- web/                  FastAPI backend + Deck.gl frontend
|-- scripts/                  запускаемые сценарии run_*.py
|-- configs/                  YAML-снимки гиперпараметров
|-- data/                     сырые и промежуточные артефакты
|-- models/                   веса обученных моделей
|-- results/                  метрики, отчёты, графики
\-- docs/                     проектная документация
\end{verbatim}
}
